\section*{NeurIPS Paper Checklist}


\begin{enumerate}

\item {\bf Claims}
    \item[] Question: Do the main claims made in the abstract and introduction accurately reflect the paper's contributions and scope?
    \item[] Answer: \answerYes{}
    \item[] Justification: The abstract and introduction state the paper's scope: verifiable agentic systems, a packed latent-mode assumption, decomposition-guided selection, and a controlled fixed-prototype experimental protocol. The discussion section separates theory-facing and deployment-facing claims.
\item {\bf Limitations}
    \item[] Question: Does the paper discuss the limitations of the work performed by the authors?
    \item[] Answer: \answerYes{}
    \item[] Justification: Section~\ref{sec:discussion} discusses the packed-mode assumption, small live-agent sample sizes, finite design menus, oracle mode labels, and current cost-measurement limitations.
\item {\bf Theory assumptions and proofs}
    \item[] Question: For each theoretical result, does the paper provide the full set of assumptions and a complete (and correct) proof?
    \item[] Answer: \answerYes{}
    \item[] Justification: Assumption~\ref{ass:packed} states the assumptions for the main theoretical result, Theorem~\ref{thm:four-term} includes a proof sketch, and Appendix~\ref{app:proofs} gives the algebraic proof and pilot concentration details.
    \item {\bf Experimental result reproducibility}
    \item[] Question: Does the paper fully disclose all the information needed to reproduce the main experimental results of the paper to the extent that it affects the main claims and/or conclusions of the paper (regardless of whether the code and data are provided or not)?
    \item[] Answer: \answerYes{}
    \item[] Justification: Section~\ref{sec:experiments} describes the AutoResearch task instantiation, task modes, model menu, router-visible signals, horizons, metrics, and evaluation protocol; Appendix~\ref{app:experiment-details} specifies the required logs, estimators, and run-metadata fields.

\item {\bf Open access to data and code}
    \item[] Question: Does the paper provide open access to the data and code, with sufficient instructions to faithfully reproduce the main experimental results, as described in supplemental material?
    \item[] Answer: \answerNo{}
    \item[] Justification: The paper is written against an existing repository and includes reproducibility commands, but an anonymized public code/data release is not included in the current submission package.

\item {\bf Experimental setting/details}
    \item[] Question: Does the paper specify all the training and test details (e.g., data splits, hyperparameters, how they were chosen, type of optimizer) necessary to understand the results?
    \item[] Answer: \answerYes{}
    \item[] Justification: Section~\ref{sec:experiments} describes the AutoResearch task instantiation, task modes, model menu, router-visible signals, primary estimators, horizons, and evaluation design. Additional implementation details are listed in Appendix~\ref{app:experiment-details}.
\item {\bf Experiment statistical significance}
    \item[] Question: Does the paper report error bars suitably and correctly defined or other appropriate information about the statistical significance of the experiments?
    \item[] Answer: \answerYes{}
    \item[] Justification: Section~\ref{sec:experiments} and Appendix~\ref{app:experiment-details} specify Wilson intervals, bootstrap intervals, survival curves, calibration diagnostics, and held-out evaluation for the reported protocol.
\item {\bf Experiments compute resources}
    \item[] Question: For each experiment, does the paper provide sufficient information on the computer resources (type of compute workers, memory, time of execution) needed to reproduce the experiments?
    \item[] Answer: \answerNo{}
    \item[] Justification: The paper reports wall-clock costs and run counts but does not provide a complete compute-resource table for every experiment.
\item {\bf Code of ethics}
    \item[] Question: Does the research conducted in the paper conform, in every respect, with the NeurIPS Code of Ethics \url{https://neurips.cc/public/EthicsGuidelines}?
    \item[] Answer: \answerYes{}
    \item[] Justification: The work studies verifiable benchmark tasks and does not involve human subjects, sensitive personal data, or released high-risk models. The submission preserves anonymity and follows NeurIPS code and data policies.

\item {\bf Broader impacts}
    \item[] Question: Does the paper discuss both potential positive societal impacts and negative societal impacts of the work performed?
    \item[] Answer: \answerYes{}
    \item[] Justification: Section~\ref{sec:discussion} notes positive uses in cost-aware and transparent agent deployment as well as limitations and failure modes. The work is methodological and does not introduce a new generative model or dataset of people.
\item {\bf Safeguards}
    \item[] Question: Does the paper describe safeguards that have been put in place for responsible release of data or models that have a high risk for misuse (e.g., pre-trained language models, image generators, or scraped datasets)?
    \item[] Answer: \answerNA{}
    \item[] Justification: The paper does not release a pre-trained model, image generator, scraped dataset, or other asset with high misuse risk. It evaluates agent harnesses on synthetic or generated benchmark tasks.
\item {\bf Licenses for existing assets}
    \item[] Question: Are the creators or original owners of assets (e.g., code, data, models), used in the paper, properly credited and are the license and terms of use explicitly mentioned and properly respected?
    \item[] Answer: \answerYes{}
    \item[] Justification: The related-work section cites existing papers and systems used for positioning, and the experimental section describes the benchmark surfaces. Released code and benchmark assets include explicit license details.
\item {\bf New assets}
    \item[] Question: Are new assets introduced in the paper well documented and is the documentation provided alongside the assets?
    \item[] Answer: \answerYes{}
    \item[] Justification: The work introduces a code harness and analysis artifacts as research assets, and Appendix~\ref{app:experiment-details} documents the required logging and configuration fields.
\item {\bf Crowdsourcing and research with human subjects}
    \item[] Question: For crowdsourcing experiments and research with human subjects, does the paper include the full text of instructions given to participants and screenshots, if applicable, as well as details about compensation (if any)?
    \item[] Answer: \answerNA{}
    \item[] Justification: The paper does not involve crowdsourcing or research with human subjects.
\item {\bf Institutional review board (IRB) approvals or equivalent for research with human subjects}
    \item[] Question: Does the paper describe potential risks incurred by study participants, whether such risks were disclosed to the subjects, and whether Institutional Review Board (IRB) approvals (or an equivalent approval/review based on the requirements of your country or institution) were obtained?
    \item[] Answer: \answerNA{}
    \item[] Justification: The paper does not involve crowdsourcing or research with human subjects, so IRB or equivalent approval is not applicable.
\item {\bf Declaration of LLM usage}
    \item[] Question: Does the paper describe the usage of LLMs if it is an important, original, or non-standard component of the core methods in this research? Note that if the LLM is used only for writing, editing, or formatting purposes and does \emph{not} impact the core methodology, scientific rigor, or originality of the research, declaration is not required.
    %this research?
    \item[] Answer: \answerYes{}
    \item[] Justification: LLMs are central to the evaluated agent designs. Section~\ref{sec:experiments} describes the three-worker menu, fixed prompt-only router, sequential signal protocol, model-identifier logging, and routing-output requirements.
\end{enumerate}
